% ===================================================================
%  TABEL Nr. 8 — Die oes. — Die jag. — Die druiweoes. — Die visvangs.
%  Meme regime que les precedents.
%
%  LE BLOC t08-c2-03-1 EST L'UN DES TROIS ECARTS DECLARES, pour une
%  raison de COMPOSITION : le fac-simile ido y ecrit « 13) » sans sa
%  parenthese ouvrante. Les traductions ecrivent le renvoi entier.
%
%  LA NOTE (*) PARLE DE L'IDO LUI-MEME et propose « mesono » pour la
%  moisson : on la traduit mot pour mot, c'est un document sur la
%  langue et non une phrase du livret.
%
%  ET LES QUATRE SAISONS FONT REVENIR LES MEMES BETES SOUS D'AUTRES
%  NOMS : les lammers du tableau 7 sont ici ooie et ramme. Le releve
%  suit leur croissance ; la colonne doit la suivre aussi.
% ===================================================================

%%K t08-noto-1 noto
\VUnotes{143.2mm}{%
\textit{(*) Omdat daardie woord nie presies is nie: gaan dit oor die vlasoes, die
druiweoes, die appeloes? Miskien sou dit goed wees om «} \VUgras{mesono} \textit{» te
gebruik wat in Mondo XIV, bl. 242 voorgestel is. Maar tot nou toe is daardie nuwe wortel
nie deur die Akademie aangeneem nie en wag dit op besprekings volgens die gebruiklike
Idistiese reglement.}%
}

%%K t08-apar-1 apar
\VUsaut{5.0mm}
\VUtitre{15.0pt}{60}{TABEL Nr. 8}
\VUsaut{3.0mm}
\VUcentre{13.2pt}{90}{\textit{Eerste tafereel.}}
\VUsaut{3.0mm}
\VUpk{10.2pt}{40}{Die oes (*).}
\VUsaut{4.0mm}
\VUpk{10.2pt}{40}{Die aansigte van die platteland.}
\VUsaut{3.0mm}

%%K t08-c1-01-1 p
1. --- Met die \VUgras{somer} het die aansig van die platteland nog 'n keer verander. Die
saad wat aan die voor toevertrou is, het gekiem; 'n groen plantjie het uit die grond
gekom en are gekry. Die graan het gevorm, daarna het dit ryp geword en van nou af hoef 'n
mens dit net te oes.

%%K t08-c1-02-1 p
2. --- Die \VUgras{veldblomme} is oop. \VUgras{Koringblomme} \textsuperscript{(1)},
\VUgras{papawers} \textsuperscript{(2)} en \VUgras{vingerhoedjies}
\textsuperscript{(3)} meng by die eenvormige skakering van die koringlande die vrolike
teenstelling van hulle vars \VUgras{kroonblare}.

%%K t08-c1-03-1 p
3. --- Die vrugtebome het hulle met blare getooi; die vrug het ryp geword. Die klein
\VUgras{kersieboom} \textsuperscript{(4)} is oorvloedig van mooi \VUgras{kersies}
\textsuperscript{(5)} wat die kinders met groot plesier pluk en eet.

%%K t08-c1-04-1 p
4. --- Van nou af kan die vee die hele dag in die buitelug deurbring.

%%K t08-c1-04-2 p
Die \VUgras{lammers} \textsuperscript{(6)} het groter geword en mooi \VUgras{ooie}
\textsuperscript{(7)} en mooi \VUgras{ramme} \textsuperscript{(8)} met 'n digte wolvag
geword. Hulle wei rustig die \VUgras{gras} van die \VUgras{weiveld}
\textsuperscript{(9)} wat met \VUgras{madeliefies} bont gemaak is.

%%K t08-c1-04-3 p
Binnekort sal 'n mens \VUgras{daardie} diere begin skeer om hulle \VUgras{wol} te kry
waarvan die laken van ons klere gemaak word.

%%K t08-tit-2 sub
\VUsaut{5.0mm}
\VUpk{10.2pt}{40}{Die herder.}
\VUsaut{3.4mm}

%%K t08-c1-05-1 p
5. --- Die \VUgras{herder} \textsuperscript{(10)} lyk of hy min op hulle let. Hy gee net
om vir die storm wat aan die horison verbygaan, en die \VUgras{skaapwagter}
\textsuperscript{(13)}, wat 'n \VUgras{draailier} \textsuperscript{(14)} na sy lippe
bring, lyk nog onbesorgder.

%%K t08-c1-06-1 p
6. --- Die hitte is drukkend; want die \VUgras{hond} \textsuperscript{(15)} kom ook sy
dors les; hy lek die vars water van die \VUgras{stroompie} \textsuperscript{(16)} en
laat 'n arme \VUgras{paddatjie} \textsuperscript{(17)} skrik wat in die oewergras
weggekruip het. Hierdie een spring in die helder water waar die \VUgras{waterlelies}
\textsuperscript{(18)} blom.

%%K t08-c1-07-1 p
7. --- 'n \VUgras{Kwikstertjie} \textsuperscript{(19)} bad by die \VUgras{fonteintjie}
\textsuperscript{(20)} wat tussen twee klippe uitspuit, en kruip weg onder die
\VUgras{doringstruike} \textsuperscript{(21)} en die \VUgras{haagdoringbosse}
\textsuperscript{(22)}.

%%K t08-tit-3 sub
\VUsaut{5.0mm}
\VUpk{10.2pt}{40}{Die oes.}
\VUsaut{3.4mm}

%%K t08-c1-08-1 p
8. --- Vir die plattelandse kinders is die koringoes 'n baie vermaaklike tyd. Hulle
\VUgras{buitel}\textsuperscript{(24)} vrolik op die \VUgras{strooihope}
\textsuperscript{(23)}, hulle help die \VUgras{oesers} \textsuperscript{(26)} en terwyl
dié die \VUgras{koringland} \textsuperscript{(27)} met hulle \VUgras{sense}
\textsuperscript{(25)} maai wat goed met die \VUgras{slypsteen} \textsuperscript{(30)}
geslyp is, versamel hulle die koring en bind dit in \VUgras{gerwe}
\textsuperscript{(31)} of in \VUgras{bondeltjies} \textsuperscript{(32)}.

%%K t08-c1-09-1 p
9. --- Soms laat 'n mens die grootstes onder hulle toe om met 'n \VUgras{sekel}
\textsuperscript{(12)} die \VUgras{goue halms} \textsuperscript{(28)} te sny wat die swaar
\VUgras{are} \textsuperscript{(29)} vol graan dra; en die kleintjies, wat die
\VUgras{vee} \textsuperscript{(33)} oppas wat in die \VUgras{weiveld}
\textsuperscript{(34)} onder die skaduwee van die groot \VUgras{lindeboom}
\textsuperscript{(35)} wei, vang met hulle \VUgras{net} \textsuperscript{(36)} die mooi
\VUgras{skoenlappers}.

%%K t08-c1-09-2 p
Sommige klim selfs op die \VUgras{waentjie} \textsuperscript{(37)} om die gerwe te
rangskik wat met 'n \VUgras{vurk} \textsuperscript{(38)} aan hulle aangegee word.

%%K t08-c1-10-1 p
10. --- Op die hele \VUgras{vlakte} \textsuperscript{(39)} waar die \VUgras{lewerikke}
\textsuperscript{(41)} opvlieg, heers die lewendigheid. Almal is met die oes besig. Maar
terwyl die armes aanhou om die ou werktuie te gebruik, \VUgras{sense, sekels} en
\VUgras{dorsvleëls} \textsuperscript{(40)}, laat die ryk eienaars hulle koring of hulle
\VUgras{rog} \textsuperscript{(43)} met 'n \VUgras{maaimasjien} \textsuperscript{(42)}
maai. Daarna, omdat hulle nie die gerwe na die skuur hoef te dra nie, maak 'n mens groot
\VUgras{gerfhope} \textsuperscript{(44)}.

%%K t08-c1-11-1 p
11. --- Dan bring 'n mens 'n \VUgras{dorsmasjien} \textsuperscript{(47)} aan wat 'n
\VUgras{trekstoommasjien} \textsuperscript{(45)} met 'n \VUgras{dryfriem}
\textsuperscript{(46)} beweeg, en so word die graan van die \VUgras{strooi} geskei sonder
om die land te verlaat waar dit gesaai is.

%%K t08-tit-4 sub
\VUsaut{5.0mm}
\VUpk{10.2pt}{40}{Die storm.}
\VUsaut{3.4mm}

%%K t08-c1-12-1 p
12. --- Dikwels gebeur \VUgras{groot storms} \textsuperscript{(53)} in die oestyd. Eers
word die gerommel van die \VUgras{donder} van ver gehoor; 'n \VUgras{stormwind uit die
weste} \textsuperscript{(54)} begin waai en buig hygend die dikste bome van die
\VUgras{bos} \textsuperscript{(49)}. Swart \VUgras{wolkmassas} \textsuperscript{(56)}
bestorm die hemel; hulle word deurkruis deur die verblindende sigsaglyne van die
\VUgras{weerlig} \textsuperscript{(55)}. Die \VUgras{reën} en soms die \VUgras{hael}
begin val en plet die oeste wat gereed is om gemaai te word.

%%K t08-c1-13-1 p
13. --- Dikwels steek die weerlig 'n gebou aan die brand. So is die dak van die
\VUgras{windmeul} \textsuperscript{(50)} verlede jaar tot as verbrand en twee van sy
\VUgras{groot vlerke} \textsuperscript{(51)} verbrysel.

%%K t08-c1-14-1 p
14. --- Na 'n paar uur word die kalmte weer gebore. 'n Blink \VUgras{reënboog}
\textsuperscript{(52)} verskyn aan die horison en die natuur hervat sy lewe wat 'n paar
oomblikke onderbreek is. Die plattelanders, wat in die \VUgras{hutte}
\textsuperscript{(48)} geskuil het, begin weer werk en probeer dapper om die skade te
herstel wat hulle pas gely het.

%%K t08-tit-5 sub
\VUsaut{6.0mm}
\VUcentre{13.2pt}{90}{\textit{Tweede tafereel.}}
\VUsaut{3.6mm}

%%K t08-tit-6 sub
\VUpk{10.2pt}{40}{Die jag. --- Die druiweoes.}
\VUsaut{1.4mm}
\VUpk{10.2pt}{40}{Die visvangs.}
\VUsaut{4.0mm}

%%K t08-c2-01-1 p
1. --- Teen die einde van \VUgras{Augustus} of in die middel van \VUgras{September},
volgens die streke, val die opening van die jagseisoen. Skaars het die eerste sonstrale
die \VUgras{akkedisse} \textsuperscript{(2)} uit hulle gat laat kom, of die
\VUgras{jagter} \textsuperscript{(5)} gaan op pad met sy \VUgras{honde}
\textsuperscript{(4)}.

%%K t08-c2-02-1 p
2. --- Hy het sy stewige \VUgras{jagpak} \textsuperscript{(9)} van dik laken aangetrek,
hom met leer \VUgras{kamaste} geskoei waarmee hy sal kan loop in die vogtige gras waar
die \VUgras{paddas} \textsuperscript{(1)} wegkruip; hy het sy \VUgras{geweer}
\textsuperscript{(6)} en sy \VUgras{wildsak} \textsuperscript{(10)} geneem, en hier
vertrek hy.

%%K t08-c2-03-1 p
3. --- Die ywer van die agtervolging bring hom in die middel van 'n wingerd waar die
druiweoes baie lewendig is. Die kinders van die wynboer, wat gewoonlik die bokke op die
pad oppas, laat hulle vandag op goeie geluk wei, en nadat hulle 'n ou \VUgras{bok}
\textsuperscript{(3)} aan 'n paal vasgemaak het, wat nie sou nalaat om sy vryheid te
gebruik om te ontsnap nie, ken hulle op hulle beurt aan hulleself 'n deel aan die plesier
toe. Die een gryp 'n tros druiwe in 'n \VUgras{oesmandjie} \textsuperscript{(12)} en
vermaak hom deur die \VUgras{sap} \textsuperscript{(14)} in 'n bekertjie
\textsuperscript{(13)} te laat loop om mos te maak (wyn wat nie gegis het nie). Die ander
kap sy hare met 'n \VUgras{blarekrans} \textsuperscript{(15)} wat hy pas gevleg het.

Naby hulle kom skaamtelose \VUgras{mossies} \textsuperscript{(16)} selfs voor die pers om
die druiwe te buit.

%%K t08-c2-04-1 p
4. --- Terwyl die kinders hulle vermaak, werk die manne. Die \VUgras{druiweplukkers}
\textsuperscript{(18)} bring die druiwe in \VUgras{rugmandjies} \textsuperscript{(17)}
in die kelder; 'n \VUgras{kuiper} \textsuperscript{(19)} maak die \VUgras{vate}
\textsuperscript{(20)} reg, gaan na of die \VUgras{duie} \textsuperscript{(22)} en die
\VUgras{bodems} \textsuperscript{(21)} in goeie toestand is, en vervang die gebreekte
\VUgras{hoepels} \textsuperscript{(23)}. Hy het ook aan die \VUgras{kuipe}
\textsuperscript{(25)} gewerk waarin 'n mens die \VUgras{druiwe} \textsuperscript{(26)}
gooi om hulle te laat gis.

%%K t08-c2-05-1 p
5. --- Wanneer die gisting klaar is, tap 'n mens die wyn af en sit die oorskiet op die
\VUgras{pers} \textsuperscript{(28)}, en deur geleidelik met die groot \VUgras{skroef}
\textsuperscript{(29)} te klem, pers 'n mens die sap uit wat in 'n \VUgras{kuipie}
\textsuperscript{(32)} loop, en nog kry 'n mens \VUgras{wyn} \textsuperscript{(31)}.

%%K t08-tit-8 sub
\VUsaut{6.0mm}
\VUpk{10.2pt}{40}{Bier en appelwyn.}
\VUsaut{3.4mm}

%%K t08-c2-06-1 p
6. --- Teen die muur van die \VUgras{kelder} \textsuperscript{(27)} klim 'n tak
\VUgras{hop} \textsuperscript{(30)} op. Daar is dit net 'n sierplant, maar in party lande,
waar die wingerdstok baie skaars is, word die hop gekweek om \VUgras{bier} te maak.

%%K t08-c2-07-1 p
7. --- Die klein \VUgras{appelboom} \textsuperscript{(33)} is gelaai met appels wat,
wanneer hulle ryp is, uitstekende \VUgras{appelwyn} sal lewer. In sy \VUgras{hok}
\textsuperscript{(34)} kyk die \VUgras{kanarie} \textsuperscript{(35)} en die
\VUgras{putter} \textsuperscript{(36)} met afgunstige oë na die mossies wat vry
rondbaljaar.

%%K t08-c2-08-1 p
8. --- So ver 'n mens kan sien, op die helling van die heuwels, staan
\VUgras{wingerdpale} \textsuperscript{(40)} in rye wat die pragtige
\VUgras{wingerdstokke} \textsuperscript{(39)} dra, gelaai met swaar \VUgras{trosse}.

%%K t08-c2-08-2 p
Die hele jaar het die \VUgras{wynboer} \textsuperscript{(42)} gewerk om sy
\VUgras{wingerd} \textsuperscript{(38)} te versorg en nou word hy vir sy moeite beloon
met 'n mooi oes. Die \VUgras{druiweplukstervroue} \textsuperscript{(41)} maak die kuipe
vol wat op die waentjie gesit is wat deur \VUgras{muile} \textsuperscript{(43)} getrek
word, en almal is vrolik.

%%K t08-tit-9 sub
\VUsaut{5.0mm}
\VUpk{10.2pt}{40}{Die visvangs.}
\VUsaut{3.4mm}

%%K t08-c2-09-1 p
9. --- Net so is dit met die \VUgras{hengelaars} \textsuperscript{(44)} wat van die
oggend af aan die oewer van die water kom sit het met hulle \VUgras{hengelstokke}
\textsuperscript{(45)}.

%%K t08-c2-09-2 p
Die \VUgras{vis} \textsuperscript{(47)} hap die \VUgras{hoek} sodra hulle hulle
\VUgras{lyn} \textsuperscript{(46)} in die water laat sak, en skaars het hulle tyd om die
\VUgras{aas} te vernuwe of 'n nuwe slagoffer spartel aan die punt van die lyn.

%%K t08-c2-09-3 p
Tog vang die \VUgras{netvisser} \textsuperscript{(48)}, wat van sy \VUgras{boot}
\textsuperscript{(49)} die \VUgras{werpnet} in die middel van die \VUgras{rivier}
\textsuperscript{(50)} gooi, meer visse.

%%K t08-c2-10-1 p
10. --- Die herfs is die seisoen van die goeie uitstappies: te voet, \VUgras{te perd}
\textsuperscript{(52)}, \VUgras{per fiets} \textsuperscript{(53)}, \VUgras{per motor}
\textsuperscript{(54)}. Die \VUgras{paaie} \textsuperscript{(51)} is skoon en 'n mens
gebruik die laaste mooi dae, want hier kom die \VUgras{swaeltjies}
\textsuperscript{(56)} reeds op die dakke bymekaar om met volle vlerke na warmer lande
te vlieg. Die blaredak van die ou \VUgras{okkerneutboom} \textsuperscript{(57)} word geel,
die \VUgras{neute} val en die hoë \VUgras{bome} wat in 'n groep soos 'n \VUgras{ruiker}
\textsuperscript{(58)} op die \VUgras{hoogte} \textsuperscript{(59)} staan, sal binnekort
hulle blare verloor.

%%K t08-c2-11-1 p
11. --- Die \VUgras{son} \textsuperscript{(60)} kom later op, die dae word korter, 'n
taamlik bytende koue begin snags gevoel word, in die \VUgras{oggend}, en hier verlaat
vlugte \VUgras{houtsnippe} \textsuperscript{(61)} reeds ons ongasvrye streke.
\VUsaut{6.0mm}
\VUfilet{22mm}
